\title{EE5311 -- Assignment 2}
\author{Joyen Benitto}
\date{\today}
\maketitle

\section{Problem 1: Static CMOS Inverter}

\subsection*{Device Parameters}

The following long-channel/velocity-saturation parameters (sky130, 1.8 V process) are used throughout this problem, as in Assignment 1:

\begin{itemize}
  \item nMOS: $V_{Tn} = 0.7$ V, $\mu_n = 0.025\ \text{m}^2/\text{Vs}$, $C_{oxn} = 8.34\ \text{fF}/\mu\text{m}^2$, $v_{satn} = 8\times10^4$ m/s, $\lambda_n = 0.2/\text{V}$
  \item pMOS: $|V_{Tp}| = 0.7$ V, $\mu_p = 0.009\ \text{m}^2/\text{Vs}$, $C_{oxp} = 8.16\ \text{fF}/\mu\text{m}^2$, $v_{satp} = 3\times10^4$ m/s, $\lambda_p = 0.2/\text{V}$
\end{itemize}

\subsection{Sizing $(W/L)_p$ for $V_{th} = V_{DD}/2$}

Consider the static CMOS inverter with the nMOS fixed at $W/L = 0.42/0.15\ \mu\text{m}$. Find $(W/L)_p$ such that the inverter threshold voltage $V_{th} = V_{DD}/2$, and plot the DC transfer characteristic (VTC).

\subsubsection{Calculation}

% TODO: derive the sizing condition for Vth = VDD/2 (equal currents at Vin=Vout=VDD/2,
% both devices in saturation) using the velocity-saturation model from Assignment 1,
% and solve for (W/L)_p given (W/L)_n = 0.42/0.15.
\begin{center}
  \fbox{\parbox{0.6\textwidth}{\centering\vspace{2.5cm}Insert $(W/L)_p$ sizing derivation here\vspace{2.5cm}}}
\end{center}

The value implemented in simulation is $W_p = 1.1995\ \mu\text{m}$, $L_p = 0.15\ \mu\text{m}$ (i.e.\ $(W/L)_p \approx 8$), against the fixed nMOS $W_n/L_n = 0.42/0.15$.

\subsubsection{Schematic}

\begin{figure}[h]
  \centering
  % TODO: replace with the actual xschem screenshot of assign1a.sch
  \fbox{\parbox{0.6\textwidth}{\centering\vspace{3cm}Insert screenshot of \texttt{assign1a.sch} here\vspace{3cm}}}
  \caption{Static CMOS inverter with $(W/L)_n = 0.42/0.15$, $(W/L)_p$ sized for $V_{th}=V_{DD}/2$.}
  \label{fig:1-schematic}
\end{figure}

\subsubsection{Measurement}

The VTC and its derivative were generated by the following ngspice control block in \texttt{assign1a.spice} (\texttt{sim\_type = 1} branch):

\begin{verbatim}
.control
  let sim_type = 1

  if sim_type = 1
    * Run DC Sweep for VTC and Noise Margins
    destroy all
    dc Vin 0 1.8 0.01
    run
    plot v(Vout) v(Vin) vs v(Vin)
    plot deriv(v(Vout))
  else
    * Run Transient Pulse Response
    destroy all
    tran 0.01ns 20ns
    run
    plot v(Vin) v(Vout) vs time
  end
.endc
\end{verbatim}

\begin{figure}[h]
  \centering
  % TODO: insert 1_vtc.png (v(Vout) vs v(Vin))
  \fbox{\parbox{0.8\textwidth}{\centering\vspace{3cm}Insert VTC plot here\vspace{3cm}}}
  \caption{Simulated DC transfer characteristic $V_{out}$ vs.\ $V_{in}$.}
  \label{fig:1-vtc}
\end{figure}

\paragraph{Discussion.} % TODO: does the simulated Vth land at VDD/2 = 0.9V? report the actual value and any discrepancy vs. the hand calculation.

\subsection{Part (a): Derivative and Noise Margins}

\subsubsection{Calculation}

The noise margins are found from the unity-gain points of the VTC, i.e.\ where $\left|\frac{dV_{out}}{dV_{in}}\right| = 1$:
\[
NM_H = V_{OH} - V_{IH}, \qquad NM_L = V_{IL} - V_{OL}
\]
where $V_{IL}$, $V_{IH}$ are the input voltages at the unity-gain points, and $V_{OL}$, $V_{OH}$ are the corresponding output voltages.

\subsubsection{Measurement}

\begin{figure}[h]
  \centering
  % TODO: insert 1a_derivative.png (deriv(v(Vout)) vs v(Vin))
  \fbox{\parbox{0.8\textwidth}{\centering\vspace{3cm}Insert $dV_{out}/dV_{in}$ plot here\vspace{3cm}}}
  \caption{Derivative of the VTC, used to locate the unity-gain points $V_{IL}$, $V_{IH}$.}
  \label{fig:1a-derivative}
\end{figure}

\begin{table}[h]
  \centering
  \begin{tabular}{|c|c|}
    \hline
    Quantity & Value (V) \\
    \hline
    $V_{IL}$ & TODO \\
    $V_{IH}$ & TODO \\
    $V_{OL}$ & TODO \\
    $V_{OH}$ & TODO \\
    $NM_L$   & TODO \\
    $NM_H$   & TODO \\
    \hline
  \end{tabular}
  \caption{Noise margins extracted from the unity-gain points of the VTC.}
  \label{tab:1a-nm}
\end{table}

\paragraph{Discussion.} % TODO

\subsection{Part (b): Effect of Resizing $(W/L)_p$ by $10\times$}

How does the VTC change if $(W/L)_p$ is increased by a factor of 10? Decreased by a factor of 10?

\subsubsection{Calculation}

% TODO: qualitative/analytical expectation -- increasing (W/L)_p strengthens the pMOS
% pull-up relative to the (fixed) nMOS pull-down, shifting Vth up toward VDD; decreasing
% (W/L)_p weakens the pMOS, shifting Vth down toward 0.

\subsubsection{Schematic}

\begin{figure}[h]
  \centering
  % TODO: replace with xschem screenshot of assign1b.sch
  \fbox{\parbox{0.6\textwidth}{\centering\vspace{3cm}Insert screenshot of \texttt{assign1b.sch} here\vspace{3cm}}}
  \caption{Inverter with pMOS $L_p = 1.5\ \mu\text{m}$ ($(W/L)_p$ decreased by $10\times$ relative to Part 1's sizing).}
  \label{fig:1b-schematic}
\end{figure}

\texttt{assign1b.spice} implements the $(W/L)_p$ decreased-by-$10\times$ case ($L_p$ scaled from $0.15\,\mu\text{m}$ to $1.5\,\mu\text{m}$, $W_p$ unchanged). % TODO: add the increased-by-10x variant (e.g. Wp x10, or Lp/10) as a second schematic/netlist.

\subsubsection{Measurement}

\begin{figure}[h]
  \centering
  % TODO: insert 1b_vtc_compare.png -- overlay original, x10, /10 VTCs
  \fbox{\parbox{0.8\textwidth}{\centering\vspace{3cm}Insert overlaid VTC comparison (original, $(W/L)_p\times10$, $(W/L)_p/10$) here\vspace{3cm}}}
  \caption{VTC comparison as $(W/L)_p$ is scaled by $10\times$ and $1/10\times$ relative to Part 1's sizing.}
  \label{fig:1b-vtc}
\end{figure}

\paragraph{Discussion.} % TODO: report how Vth shifts in each direction and tie back to the pull-up/pull-down strength argument above.

\subsection{Part (c): VTC Across Supply Voltages (0.2 V to 1.8 V)}

Obtain the DC transfer characteristic for $V_{DD}$ ranging from 0.2 V to 1.8 V in steps of 0.2 V. Plot $I_{DS}$ vs.\ $V_{in}$ for all these supply voltages. Can it be used as an inverter for these supply voltages? If not, what change to the pMOS width or length would fix it?

\subsubsection{Calculation}

% TODO: at low VDD, both devices' |VDS|, |VGS| may fall below their thresholds (Vtn, |Vtp| = 0.7V),
% so neither device may turn on -- the circuit fails to behave as an inverter below some VDD.
% State the expected cutoff VDD and the proposed pMOS resizing fix.

\subsubsection{Schematic}

\begin{figure}[h]
  \centering
  % TODO: replace with xschem screenshot of assing1c.sch
  \fbox{\parbox{0.6\textwidth}{\centering\vspace{3cm}Insert screenshot of \texttt{assing1c.sch} here\vspace{3cm}}}
  \caption{Static inverter with $V_{DD}$ (\texttt{Vin1}) parametrized as a second DC sweep source.}
  \label{fig:1c-schematic}
\end{figure}

\subsubsection{Measurement}

The family of VTCs across $V_{DD}$ was generated by a nested DC sweep in \texttt{assing1c.spice}:

\begin{verbatim}
* Nested DC Sweep: Sweep Vin from 0 to 1.8, and loop Vin1 (VDD) from 0.2 to 1.8
.dc Vin 0 1.8 0.01 Vin1 0.2 1.8 0.2

.control
run
plot v(Vout) vs v(Vin)
.endc
\end{verbatim}

\begin{figure}[h]
  \centering
  % TODO: insert 1c_vtc_family.png
  \fbox{\parbox{0.8\textwidth}{\centering\vspace{3cm}Insert $V_{out}$ vs.\ $V_{in}$ family (one curve per $V_{DD}$) here\vspace{3cm}}}
  \caption{VTC family for $V_{DD} = 0.2,0.4,\ldots,1.8$~V.}
  \label{fig:1c-vtc-family}
\end{figure}

\begin{figure}[h]
  \centering
  % TODO: add an Ids vs Vin plot/measurement for the same VDD sweep -- not yet in assing1c.spice
  \fbox{\parbox{0.8\textwidth}{\centering\vspace{3cm}Insert $I_{DS}$ vs.\ $V_{in}$ family (one curve per $V_{DD}$) here\vspace{3cm}}}
  \caption{$I_{DS}$ vs.\ $V_{in}$ family for $V_{DD} = 0.2,0.4,\ldots,1.8$~V.}
  \label{fig:1c-ids-family}
\end{figure}

\paragraph{Discussion.} % TODO: identify which VDD values fail to produce inverter-like (full-swing,
% monotonic, sharp-transition) behavior, and state the proposed pMOS width/length fix.

\subsection{Part (d): Peak $I_{DS}$ at $V_{DD} = 0.8, 1.8$~V vs.\ Analytical}

For $V_{DD} = 0.8$ V and $1.8$ V, find the peak $I_{DS}$ and compare it with the analytical value.

\subsubsection{Calculation}

% TODO: peak IDS occurs near Vin = Vth, where both transistors are in saturation;
% use the velocity-saturation saturation-current expression from Assignment 1 for each device
% and solve self-consistently (or evaluate at the simulated Vth) for the analytical peak current.

\subsubsection{Measurement}

% TODO: extract peak IDS from the VDD=0.8V and VDD=1.8V curves (e.g. via `meas dc` in ngspice,
% or from the 1c sweep data / a dedicated sim1d.cir).

\begin{table}[h]
  \centering
  \begin{tabular}{|c|c|c|}
    \hline
    $V_{DD}$ (V) & Simulated peak $I_{DS}$ & Analytical peak $I_{DS}$ \\
    \hline
    0.8 & TODO & TODO \\
    1.8 & TODO & TODO \\
    \hline
  \end{tabular}
  \caption{Peak $I_{DS}$: simulated vs.\ analytical, for $V_{DD}=0.8$ and $1.8$~V.}
  \label{tab:1d-peak-ids}
\end{table}

\paragraph{Discussion.} % TODO

\section{Problem 2: Pseudo-nMOS Inverter}

Consider the pseudo-nMOS inverter: the pMOS gate is tied permanently to GND (always on), and the nMOS ($W/L = 0.42/0.15\ \mu\text{m}$) is driven by $V_{in}$ as the switching device.

\subsection{Sizing the pMOS for $V_{OL} = 0.1$~V}

Size the pMOS transistor so that $V_{OL} = 0.1$ V. Using Ngspice, plot the DC transfer characteristic.

\subsubsection{Calculation}

% TODO: at Vin = VDD (Vout = VOL = 0.1V), the nMOS is deep in triode and the pMOS is in
% saturation (Vsg = VDD, Vsd = VDD - VOL is large); equate the nMOS triode current to the
% pMOS saturation current at this operating point and solve for (W/L)_p.
\begin{center}
  \fbox{\parbox{0.6\textwidth}{\centering\vspace{2.5cm}Insert $(W/L)_p$ sizing derivation for $V_{OL}=0.1$~V here\vspace{2.5cm}}}
\end{center}

The value implemented in simulation is $W_p = 0.42\ \mu\text{m}$, $L_p = 0.30\ \mu\text{m}$ (i.e.\ $(W/L)_p$ halved relative to the nMOS), pMOS gate tied to GND.

\subsubsection{Schematic}

\begin{figure}[h]
  \centering
  % TODO: replace with xschem screenshot of assign2a.sch
  \fbox{\parbox{0.6\textwidth}{\centering\vspace{3cm}Insert screenshot of \texttt{assign2a.sch} here\vspace{3cm}}}
  \caption{Pseudo-nMOS inverter: pMOS gate tied to GND (always on), pMOS sized for $V_{OL}=0.1$~V.}
  \label{fig:2-schematic}
\end{figure}

\subsubsection{Measurement}

\begin{verbatim}
.control
dc Vin 0.1 1.8 0.01
plot v(Vout) vs v(Vin)
.endc
\end{verbatim}

\begin{figure}[h]
  \centering
  % TODO: insert 2_vtc.png
  \fbox{\parbox{0.8\textwidth}{\centering\vspace{3cm}Insert pseudo-nMOS VTC plot here\vspace{3cm}}}
  \caption{Simulated DC transfer characteristic of the pseudo-nMOS inverter.}
  \label{fig:2-vtc}
\end{figure}

\paragraph{Discussion.} % TODO: report the simulated VOL and compare against the 0.1V target.

\subsection{Part (a): Inverter Threshold Voltage}

Find the inverter threshold voltage $V_{th}$ (where $V_{in} = V_{out}$ on the VTC of Section~2.1).

\subsubsection{Measurement}

% TODO: read Vth off the same VTC data (fig:2-vtc) at the point v(Vin) = v(Vout).

\paragraph{Discussion.} % TODO

\subsection{Part (b): Noise Margins}

Plot the derivative and obtain the high and low noise margins. Do an approximate calculation using the analytical models and compare with the simulated value.

\subsubsection{Calculation}

% TODO: analytical NMH/NML for the pseudo-nMOS inverter, using the same
% unity-gain-point definition as Problem 1(a) but with the always-on pMOS load.

\subsubsection{Measurement}

The derivative was generated from \texttt{assign2b.spice}:

\begin{verbatim}
.control
dc Vin 0 1.8 0.01
plot deriv(v(Vout))
print -i(Vin1) * 1.8
.endc
\end{verbatim}

\begin{figure}[h]
  \centering
  % TODO: insert 2b_derivative.png
  \fbox{\parbox{0.8\textwidth}{\centering\vspace{3cm}Insert $dV_{out}/dV_{in}$ plot here\vspace{3cm}}}
  \caption{Derivative of the pseudo-nMOS VTC, used to locate $V_{IL}$, $V_{IH}$.}
  \label{fig:2b-derivative}
\end{figure}

\begin{table}[h]
  \centering
  \begin{tabular}{|c|c|c|}
    \hline
    Quantity & Analytical & Simulated \\
    \hline
    $NM_L$ & TODO & TODO \\
    $NM_H$ & TODO & TODO \\
    \hline
  \end{tabular}
  \caption{Noise margins: analytical vs.\ simulated.}
  \label{tab:2b-nm}
\end{table}

\paragraph{Discussion.} % TODO

\subsection{Part (c): Average Power Dissipated at $V_{in} = V_{DD}$}

Assuming $V_{in} = V_{DD}$, find the average power dissipated and compare with the analytically obtained value.

\subsubsection{Calculation}

% TODO: at Vin = VDD, Vout = VOL and the static (pull-through) current I_stat flows continuously
% from VDD to GND through the always-on pMOS and the fully-on nMOS; average power
% P_avg = VDD * I_stat (no switching component, since Vin is held static). Evaluate I_stat
% analytically at this operating point using the same device equations as Section 2.1.

\subsubsection{Measurement}

\texttt{assign2b.spice} prints \texttt{-i(Vin1) * 1.8} across the sweep; the value at $V_{in} = V_{DD} = 1.8$~V gives the static power at this operating point. % TODO: report the printed value at Vin=1.8V specifically.

\begin{table}[h]
  \centering
  \begin{tabular}{|c|c|}
    \hline
    Quantity & Value \\
    \hline
    Analytical $P_{avg}$ & TODO \\
    Simulated $P_{avg}$  & TODO \\
    \hline
  \end{tabular}
  \caption{Average power dissipated at $V_{in}=V_{DD}$: analytical vs.\ simulated.}
  \label{tab:2c-power}
\end{table}

\paragraph{Discussion.} % TODO
